% !TeX root = ../papers/nc-p3.tex
\section{印付き四平面対の交換子階数と階数一の差}
\label{sec:paired-marked-four-planes}
定理\ref{thm:marked-point-commutators}の印で $g=q=4$ とする。
$D=\langle G,F\rangle$, $D'=\langle Q,P\rangle$ とし、$[D,D]$ は交換子の線型包を表す。
\begin{theorem}\label{thm:paired-marked-four-planes}
各三列が独立なら $r=\dim[D,D]=\dim[D',D']$ は4か5である。
両側で同時に基底を選び
\[
 F_x=G_x+\alpha(x)T,\qquad P_y=Q_y+\beta(y)L
\]
とできる。ここで $\alpha,\beta$ は非零一次形式で、$T,L$ は各三列の包の外にある。
このとき $[G_{\alpha^t},T]=[Q_{\beta^t},L]=0$ である。
$r=5$ では各四平面に可換二平面がただ一つある。
$r=4$ では各四平面の交換子核は二次元で、表示
$\ker(\wedge^2D\to[D,D])=h\wedge E$ を持つ。ここで $E$ は三次元で
$h\in E$ はスカラーを加えれば階数1であり、可換二平面は $k h$ を共有する射影直線をなす。
\end{theorem}
\begin{proof}
$G,F$ の二つの像超平面の交わりを
$G_{a(t)}=F_{c(t)}$ と書く。$a,c$ は二次元の $t$ に関する単射である。
$G_a\wedge F_c$ は $\wedge^2D$ を張る（交叉の二基底と各補元の対を取る）。
従って両側の交換子次元は交換子テンソルの行列階数 $r$ に一致する。
印付き等式から
\[
 \mathcal S=\langle c(t)a(t)^t\rangle\subset [D',D']^\perp
\]
が跡対合に関し成立する。$\mathcal S$ は二次元から二次元への対称二行列全体を
単射と全射で挟んだ三次元空間で、全体で共通核を持つ。単位行列も直交するので $r\le5$。
逆側にも同じ構成がある。

どちらの四平面も非退化対称形式について自己随伴にはならない。
例えば $D'$ を対称にできれば印付き等式より全ての $x$ で $[G_x,F_x]=0$ である。
$D$ の座標を交叉二元と補元二つで選び
$G_x=(u_1,u_2,a,0)$, $F_x=(v_1,v_2,0,b)$ と書く（文字は $x$ の一次式）。
$x\mapsto G_x\wedge F_x$ の座標二次式は
$a\langle v_1,v_2,b\rangle+b\langle u_1,u_2,a\rangle$
を含む。各括弧は全一次式空間なので $a,b$ が独立ならその次元は5となる。
一方この二次写像の像の包は補題\ref{lem:irreducible-four-plane-brackets}により
次元3以下の交換子核に入る。よって $b$ は $a$ の倍、全座標二次式は $a$ で割り切れる。
$\ker a$ 上で $u_1v_2-u_2v_1=0$ となり、この二平面から交叉への二つの同型は比例する。
従って $F_x=dG_x+a(x)Z$、可換性により $Z$ は既約な $G$ と交換する跡零元なので零となり矛盾する。

次に、既約な跡零四平面で $r=3$ なら必ず自己随伴にできることを示す。
前掲補題の証明により、一般 $X\in D$ は三固有値相異なり、
$\dim[X,D]=3$ である（一般で2以下なら五次元の既約Jordan部分空間が生じる矛盾）。
$X$ を対角化すると $[D,D]=[X,D]$ は非対角である。
$Y$ の成分を $l_{ij}(Y)$ とすると、全ての対角交換子成分が零なので
$l_{12}\wedge l_{21}=l_{23}\wedge l_{32}=l_{31}\wedge l_{13}$。
この交代形式が非零ならその台二平面に全非対角成分が属して $[X,D]$ の階数に反する。
よって各対の二成分は比例し、三対の寄与は各一次元、互いに独立である。
三対を別々に支える元を選び、向き $12,23,31$ でその成分を $p_i$、逆成分を $q_i$ とする。
最初の二元の交換子の $13,31$ は $p_1p_2,-q_1q_2$。
$[X,D]$ の同成分対が $(q_3,-p_3)$ の倍なので $p_1p_2p_3=q_1q_2q_3$。
既約性により有向成分グラフは強連結である。一方向だけの辺があればそれを含む有向三角形があり、
この積等式で逆向き三角形も存在して矛盾する。従って全成分は非零で、
同じ積等式から $h_iY_{ij}=h_jY_{ji}$ を満たす非零対角計量が存在する。これで $r=3$ を排除した。

残る交換子核 $K\subset\wedge^2D$ の全元が可分解であることを示す。
$\phi(X,Y,Z)=\operatorname{tr}(X[Y,Z])$ と置く。
$\phi\ne0$ なら四次元の三形式には根基直線 $kh$ があり、その二重縮約の核は $h\wedge D$。
$K$ はこの核に含まれる。$\phi=0$ なら $kI+D\subset[D,D]^\perp$ より $r=4$ で等号となる。
そこで反対側からの三平面 $\mathcal S'=\{C S A:S=S^t\}$ が $kI+D$ に入る。
$C$ は $3\times2$ の単射、$A$ は $2\times3$ の全射とし、$J=AC$ は階数1以上である。
$\mathcal S'$ の三基底の三形式は符号を除き
$(J_{12}-J_{21})\det J$ となる（対称二行列の標準三基底で展開する）。
$J$ が非対称なら階数1であり、合同変換で $E_{12}$ にできる。
$C$ の二列を基底の初めに取り、第三列を選べば $\mathcal S'$ の基底は
$E_{12},E_{13}+E_{22},E_{23}$。その交換子は厳密上三角全体を張るので、
$\phi=0$ から $D$ は上三角に含まれてしまう。従って $J$ は対称である。
このとき $[CSA,CTA]=C(SJT-TJS)A$ の中項は交代二行列で、全交換子は一次元以下である。
$\mathcal S'$ を跡零化した三平面を $E$ とすると $\wedge^2 E\cap K$ は二次元以上、
$\dim K=2$ だから等号である。三次元の全二ベクトルは可分解であり、
二次元の束は共通直線 $kh\subset E$ を持つ。
$\phi\ne0$ で $r=4$ の場合にも $K\subset h\wedge D$ より同じ表示を得る。
いずれでも $h$ は $E$ 全体と交換するので三行の巡回元にはなりえず、
非零跡零であることからあるスカラー行列との差が階数1である。
$r=5$ のときは $K$ が一本の可分解直線である。

最後に $D$ への全射 $\pi:(a,c)\mapsto G_a+F_c$ を考える。
入力六空間に標準交代形式 $a^tc'-a'^tc$ を入れると、逆二ベクトルを商へ移したものは
符号を除き $\Omega=\sum_i G_i\wedge F_i$ である。交換子式の跡から $\Omega\in K$。
商の二ベクトルの核次元は $\dim(\ker\pi\cap(\ker\pi)^\perp)$ である
（商の双対は交代形式で $(\ker\pi)^\perp$ と同一視される）。
よって可分解性から二次元の $\ker\pi$ は等方的である。
$\ker\pi=\{(a(t),-c(t))\}$ なので $a(t)^tc(u)=a(u)^tc(t)$。
対応 $a(t)\mapsto c(t)$ は可逆対称行列へ延長できる。実際、指定二列は独立であり、
その主二行二列が可逆なら最後の対角値を一般に選べばよく、階数1なら最後の列との対合が
その根基で非零であるため既に行列式が非零である。合同変換でこの延長を単位行列にする。
すると $G,F$ は同じ二平面上で等しく、主張の差表示が従う。
反対側の同じ基底選択はこの一致を同時共役するだけなので両側で成立する。
交換子の跡縮約をもう一度用いると最後の二可換条件も従う。
\end{proof}
